\rSec2[expected.void]{Partial specialization of \tcode{expected} for \tcode{void} types}

\rSec3[expected.void.general]{General}

\begin{codeblock}
namespace std {
  template<class E>
  class @\libglobal{expected}@<void, E> {
  private:
    using @\exposid{error-value-type}@ = remove_cv_t<remove_reference_t<E>>; // \expos

  public:
    using @\libmember{value_type}{expected}@ = void;
    using @\libmember{error_type}{expected}@ = E;
    using @\libmember{unexpected_type}{expected}@ = unexpected<E>;

    template<class U> using rebind = expected<U, error_type>;

    // \ref{expected.void.cons}, constructors

    constexpr expected() noexcept;

    constexpr expected(const expected& rhs) noexcept(
        is_nothrow_copy_constructible_v<E>);

    constexpr expected(expected&& rhs) noexcept(is_nothrow_move_constructible_v<E>);

    template<class U, class G>
    constexpr explicit(!is_convertible_v<const G&, E>)
        expected(const expected<U, G>& rhs);

    template<class U, class G>
    constexpr explicit(!is_convertible_v<G, E>) expected(expected<U, G>&& rhs);

    template<class G>
    constexpr explicit(!is_convertible_v<const G&, E>) expected(const unexpected<G>& e);

    template<class G>
    constexpr explicit(!is_convertible_v<G, E>) expected(unexpected<G>&& e);

    template<class G>
    constexpr explicit(!is_convertible_v<G, E>)
        expected(const unexpected<G>& e) noexcept;

    template<class G>
    constexpr explicit(!is_convertible_v<G, E>) expected(unexpected<G>&& e) noexcept;

    template<class G> constexpr expected(const unexpected<G>&) = delete;

    template<class G>
      requires(is_reference_v<E> && !is_reference_v<G>)
    constexpr expected(unexpected<G>&&) = delete;

    constexpr explicit expected(in_place_t) noexcept;

    template<class... Args> constexpr explicit expected(unexpect_t, Args&&... args);

    template<class... Args> constexpr expected(unexpect_t, Args&&...) = delete;

    template<class... Args>
      requires(is_reference_v<E> && !is_constructible_v<E, Args...> &&
               !@\exposid{unexpect-dangles-v}@<E, Args...>)
    constexpr expected(unexpect_t, Args&&...) = delete;

    template<class U, class... Args>
    constexpr explicit expected(unexpect_t, initializer_list<U> il, Args&&... args);

    template<class U, class... Args>
    constexpr expected(unexpect_t, initializer_list<U>, Args&&...) = delete;

    template<class G>
    constexpr explicit(!is_convertible_v<G&, E>)
        expected(const expected<void, G&>& rhs);

    template<class G>
    constexpr explicit(!is_convertible_v<G&, E>) expected(expected<void, G&>&& rhs);

    // \ref{expected.void.dtor}, destructor

    constexpr ~expected();

    // \ref{expected.void.assign}, assignment

    constexpr expected& operator=(const expected& rhs) noexcept(
        is_nothrow_copy_constructible_v<E> && is_nothrow_copy_assignable_v<E>);

    constexpr expected& operator=(expected&& rhs) noexcept(
        is_nothrow_move_constructible_v<E> && is_nothrow_move_assignable_v<E>);

    template<class G> constexpr expected& operator=(const unexpected<G>& e);

    template<class G> constexpr expected& operator=(unexpected<G>&& e);

    template<class G> constexpr expected& operator=(const unexpected<G>& e);

    template<class G> constexpr expected& operator=(unexpected<G>&& e);

    template<class G> constexpr expected& operator=(const unexpected<G>&) = delete;

    template<class G>
      requires(is_reference_v<E> && !is_reference_v<G>)
    constexpr expected& operator=(unexpected<G>&&) = delete;

    constexpr void emplace() noexcept;

    // \ref{expected.void.swap}, swap

    constexpr void swap(expected& rhs) noexcept(is_nothrow_move_constructible_v<E> &&
                                                (is_reference_v<E> ||
                                                 is_nothrow_swappable_v<E>));

    friend constexpr void swap(expected& x, expected& y) noexcept(noexcept(x.swap(y)));

    // \ref{expected.void.obs}, observers

    constexpr explicit operator bool() const noexcept;
    constexpr bool has_value() const noexcept;

    constexpr void operator*() const noexcept;

    constexpr void value() const&;
    constexpr void value() &&;

    constexpr const E& error() const& noexcept;
    constexpr E& error() & noexcept;
    constexpr const E&& error() const&& noexcept;
    constexpr E&& error() && noexcept;

    template<class G = @\exposid{error-value-type}@>
    constexpr @\exposid{error-value-type}@ error_or(G&& def) const&;

    template<class G = @\exposid{error-value-type}@>
    constexpr @\exposid{error-value-type}@ error_or(G&& def) &&;

    template<class U> constexpr void value_or(U&&) const = delete;

    // \ref{expected.void.monadic}, monadic operations

    template<class F> constexpr auto and_then(F&& f) &;
    template<class F> constexpr auto and_then(F&& f) &&;
    template<class F> constexpr auto and_then(F&& f) const&;
    template<class F> constexpr auto and_then(F&& f) const&&;

    template<class F> constexpr auto or_else(F&& f) &;
    template<class F> constexpr auto or_else(F&& f) &&;
    template<class F> constexpr auto or_else(F&& f) const&;
    template<class F> constexpr auto or_else(F&& f) const&&;

    template<class F> constexpr auto transform(F&& f) &;
    template<class F> constexpr auto transform(F&& f) &&;
    template<class F> constexpr auto transform(F&& f) const&;
    template<class F> constexpr auto transform(F&& f) const&&;

    template<class F> constexpr auto transform_error(F&& f) &;
    template<class F> constexpr auto transform_error(F&& f) &&;
    template<class F> constexpr auto transform_error(F&& f) const&;
    template<class F> constexpr auto transform_error(F&& f) const&&;

    // \ref{expected.void.eq}, equality operators

    template<class T2, class E2>
    friend constexpr bool operator==(const expected& x, const expected<T2, E2>& y);

    template<class E2>
    friend constexpr bool operator==(const expected& x, const unexpected<E2>& e);

  private:
    bool @\exposid{has_val}@; // \expos
    union {
      unexpected<E> @\exposid{unex}@; // \expos
    };
  };
}
\end{codeblock}

\pnum
\mandates
A program that instantiates the definition of \tcode{expected<T, E>} with an \tcode{E} that is not a valid template argument for \tcode{unexpected} is ill-formed.

\pnum
\remarks
Any object of type \tcode{expected<T, E>} either represents a value of type \tcode{T}, or contains a value of type \tcode{E} nested within\iref{intro.object} it. Member \exposid{has_val} indicates whether the \tcode{expected<T, E>} object represents a value of type \tcode{T}. When \tcode{has_value()} is \tcode{false}, the error is \tcode{\exposid{unex}.error()}.

\rSec3[expected.void.cons]{Constructors}

\indexlibraryctor{expected}%
\begin{itemdecl}
constexpr expected() noexcept;
\end{itemdecl}

\begin{itemdescr}
\pnum
\ensures
\tcode{has_value()} is \tcode{true}.
\end{itemdescr}

\indexlibraryctor{expected}%
\begin{itemdecl}
constexpr expected(const expected& rhs) noexcept(is_nothrow_copy_constructible_v<E>);
\end{itemdecl}

\begin{itemdescr}
\pnum
\constraints
\tcode{is_copy_constructible_v<E>} is \tcode{true} and \tcode{is_trivially_copy_constructible_v<E>} is \tcode{false}.

\pnum
\effects
If \tcode{rhs.has_value()} is \tcode{false}, direct-non-list-initializes \exposid{unex} with \tcode{rhs.error()}.

\pnum
\ensures
\tcode{rhs.has_value() == this->has_value()}.

\pnum
\throws
Any exception thrown by the initialization of \exposid{unex}.

\pnum
\remarks
This constructor is defined as deleted unless \tcode{is_copy_constructible_v<E>} is \tcode{true} or \tcode{is_reference_v<E>} is \tcode{true}. This constructor is trivial if \tcode{is_trivially_copy_constructible_v<E>} is \tcode{true}.
\end{itemdescr}

\indexlibraryctor{expected}%
\begin{itemdecl}
constexpr expected(expected&& rhs) noexcept(is_nothrow_move_constructible_v<E>);
\end{itemdecl}

\begin{itemdescr}
\pnum
\constraints
\tcode{is_move_constructible_v<E>} is \tcode{true} and \tcode{is_trivially_move_constructible_v<E>} is \tcode{false}.

\pnum
\effects
If \tcode{rhs.has_value()} is \tcode{false}, direct-non-list-initializes \exposid{unex} with \tcode{std::move(rhs.error())}.

\pnum
\ensures
\tcode{rhs.has_value()} is unchanged; \tcode{rhs.has_value() == this->has_value()} is \tcode{true}.

\pnum
\throws
Any exception thrown by the initialization of \exposid{unex}.

\pnum
\remarks
This constructor is trivial if \tcode{is_trivially_move_constructible_v<E>} is \tcode{true}.
\end{itemdescr}

\indexlibraryctor{expected}%
\begin{itemdecl}
template<class U, class G>
  constexpr explicit(!is_convertible_v<const G&, E>) expected(const expected<U, G>& rhs);
template<class U, class G>
  constexpr explicit(!is_convertible_v<G, E>) expected(expected<U, G>&& rhs);
\end{itemdecl}

\begin{itemdescr}
\pnum
\constraints
\tcode{is_void_v<U>} is \tcode{true}; and \tcode{is_constructible_v<E, const G\&>} is \tcode{true}; and \tcode{is_constructible_v<unexpected<E>, expected<U, G>\&>} is \tcode{false}; and \tcode{is_constructible_v<unexpected<E>, expected<U, G>>} is \tcode{false}; and \tcode{is_constructible_v<unexpected<E>, const expected<U, G>\&>} is \tcode{false}; and \tcode{is_constructible_v<unexpected<E>, const expected<U, G>>} is \tcode{false}.

\pnum
\effects
If \tcode{rhs.has_value()} is \tcode{false}, direct-non-list-initializes \exposid{unex} with \tcode{rhs.error()}.

\pnum
\ensures
\tcode{rhs.has_value()} is unchanged; \tcode{rhs.has_value() == this->has_value()} is \tcode{true}.

\pnum
\throws
Any exception thrown by the initialization of \exposid{unex}.
\end{itemdescr}

\indexlibraryctor{expected}%
\begin{itemdecl}
template<class G>
  constexpr explicit(!is_convertible_v<const G&, E>) expected(const unexpected<G>& e);
template<class G>
  constexpr explicit(!is_convertible_v<G, E>) expected(unexpected<G>&& e);
\end{itemdecl}

\begin{itemdescr}
\pnum
\constraints
\tcode{is_constructible_v<E, const G\&>} is \tcode{true}.

\pnum
\effects
Direct-non-list-initializes \exposid{unex} with \tcode{e.error()}.

\pnum
\ensures
\tcode{has_value()} is \tcode{false}.

\pnum
\throws
Any exception thrown by the initialization of \exposid{unex}.
\end{itemdescr}

\indexlibraryctor{expected}%
\begin{itemdecl}
template<class G>
  constexpr explicit(!is_convertible_v<G, E>) expected(const unexpected<G>& e) noexcept;
template<class G>
  constexpr explicit(!is_convertible_v<G, E>) expected(unexpected<G>&& e) noexcept;
\end{itemdecl}

\begin{itemdescr}
\pnum
\constraints
\tcode{is_reference_v<G>} is \tcode{true}; and \tcode{is_constructible_v<E, G>} is \tcode{true}; and \tcode{reference_constructs_from_temporary_v<E, G>} is \tcode{false}.

\pnum
\effects
Initializes \exposid{unex} with \tcode{e.error()}.

\pnum
\ensures
\tcode{has_value()} is \tcode{false}.

\pnum
\remarks
This constructor never throws: the referent is bound, not copied.
\end{itemdescr}

\indexlibraryctor{expected}%
\begin{itemdecl}
template<class G>
  constexpr expected(const unexpected<G>&) = delete;
\end{itemdecl}

\begin{itemdescr}
\pnum
\constraints
\tcode{is_reference_v<E>} is \tcode{true} and \tcode{is_reference_v<G>} is \tcode{false}.

\pnum
\remarks
When \tcode{E} is a reference type, an overload taking \tcode{unexpected<G>} for a non-reference \tcode{G} is defined as deleted: the referent would live inside the (possibly temporary) source \tcode{unexpected<G>} object, and binding \tcode{E\&} to it would dangle. Use an \tcode{unexpected<E\&>} holding an external object instead.
\end{itemdescr}

\indexlibraryctor{expected}%
\begin{itemdecl}
constexpr explicit expected(in_place_t) noexcept;
\end{itemdecl}

\begin{itemdescr}
\pnum
\ensures
\tcode{has_value()} is \tcode{true}.
\end{itemdescr}

\indexlibraryctor{expected}%
\begin{itemdecl}
template<class... Args>
  constexpr explicit expected(unexpect_t, Args&&... args);
\end{itemdecl}

\begin{itemdescr}
\pnum
\constraints
\tcode{is_constructible_v<E, Args...>} is \tcode{true}.

\pnum
\effects
Direct-non-list-initializes \exposid{unex} with \tcode{std::forward<Args>(args)...}.

\pnum
\ensures
\tcode{has_value()} is \tcode{false}.

\pnum
\throws
Any exception thrown by the initialization of \exposid{unex}.
\end{itemdescr}

\indexlibraryctor{expected}%
\begin{itemdecl}
template<class... Args>
  constexpr expected(unexpect_t, Args&&...) = delete;
\end{itemdecl}

\begin{itemdescr}
\pnum
\constraints
\tcode{\exposid{unexpect-dangles-v}<E, Args...>} is \tcode{true}.

\pnum
\remarks
When \tcode{E} is a reference type, an overload with the same parameter types is defined as deleted if the single argument would bind \tcode{E\&} to a temporary, or if it is otherwise not usable to construct \tcode{E}.
\end{itemdescr}

\indexlibraryctor{expected}%
\begin{itemdecl}
template<class U, class... Args>
  constexpr explicit expected(unexpect_t, initializer_list<U> il, Args&&... args);
\end{itemdecl}

\begin{itemdescr}
\pnum
\constraints
\tcode{is_reference_v<E>} is \tcode{false}, and \tcode{is_constructible_v<E, initializer_list<U>\&, Args...>} is \tcode{true}.

\pnum
\effects
Direct-non-list-initializes \exposid{unex} with \tcode{il, std::forward<Args>(args)...}.

\pnum
\ensures
\tcode{has_value()} is \tcode{false}.

\pnum
\throws
Any exception thrown by the initialization of \exposid{unex}.

\pnum
\remarks
An overload with the same parameter types is defined as deleted when \tcode{E} is an lvalue reference type. An initializer list cannot provide the required long-lived error referent.
\end{itemdescr}

\indexlibraryctor{expected}%
\begin{itemdecl}
template<class U, class... Args>
  constexpr expected(unexpect_t, initializer_list<U>, Args&&...) = delete;
\end{itemdecl}

\begin{itemdescr}
\pnum
\constraints
\tcode{is_reference_v<E>} is \tcode{true}.

\pnum
\remarks
An overload with the same parameter types is defined as deleted when \tcode{E} is an lvalue reference type. An initializer list cannot provide the required long-lived error referent.
\end{itemdescr}

\indexlibraryctor{expected}%
\begin{itemdecl}
template<class G>
  constexpr explicit(!is_convertible_v<G&, E>) expected(const expected<void, G&>& rhs);
template<class G>
  constexpr explicit(!is_convertible_v<G&, E>) expected(expected<void, G&>&& rhs);
\end{itemdecl}

\begin{itemdescr}
\pnum
\constraints
\tcode{is_convertible_v<G\&, E>} is \tcode{true}; and \tcode{reference_constructs_from_temporary_v<E, G\&>} is \tcode{false}.

\pnum
\effects
If \tcode{rhs.has_value()} is \tcode{false}, direct-non-list-initializes \exposid{unex} with \tcode{rhs.error()}.

\pnum
\ensures
\tcode{rhs.has_value()} is unchanged; \tcode{rhs.has_value() == this->has_value()} is \tcode{true}.

\pnum
\remarks
This constructor never throws: the referent is bound, not copied. It participates in overload resolution only when \tcode{E} is a reference type, mirroring the \tcode{unexpected<G>} reference-\tcode{E} path above.
\end{itemdescr}

\rSec3[expected.void.dtor]{Destructor}

\indexlibrarydtor{expected}%
\begin{itemdecl}
constexpr ~expected();
\end{itemdecl}

\begin{itemdescr}
\pnum
\effects
If \tcode{has_value()} is \tcode{false}, destroys \exposid{unex}.

\pnum
\remarks
If \tcode{is_trivially_destructible_v<E>} is \tcode{true}, then this destructor is a trivial destructor.
\end{itemdescr}

\rSec3[expected.void.assign]{Assignment}

\indexlibrarymember{operator=}{expected}%
\begin{itemdecl}
constexpr expected& operator=(const expected& rhs) noexcept(is_nothrow_copy_constructible_v<E> &&
                                                            is_nothrow_copy_assignable_v<E>);
\end{itemdecl}

\begin{itemdescr}
\pnum
\constraints
\tcode{is_reference_v<E> || (is_copy_constructible_v<E> \&\& is_copy_assignable_v<E>)} is \tcode{true} and \tcode{(is_trivially_copy_constructible_v<E> \&\& is_trivially_copy_assignable_v<E> \&\& is_trivially_destructible_v<E>)} is \tcode{false}.

\pnum
\effects
If \tcode{this->has_value() \&\& rhs.has_value()}, no effects. Otherwise, if \tcode{this->has_value()}, equivalent to: \tcode{construct_at(addressof(\exposid{unex}), rhs.\exposid{unex}); \exposid{has_val} = false;} Otherwise, if \tcode{rhs.has_value()}, destroys \exposid{unex} and sets \exposid{has_val} to \tcode{true}. Otherwise, equivalent to \tcode{\exposid{unex} = rhs.\exposid{unex}}.

\pnum
\returns
\tcode{*this}.

\pnum
\remarks
This operator is defined as deleted unless \tcode{is_copy_assignable_v<E>} is \tcode{true} or \tcode{is_reference_v<E>} is \tcode{true} and \tcode{is_copy_constructible_v<E>} is \tcode{true} or \tcode{is_reference_v<E>} is \tcode{true}. This operator is trivial if \tcode{is_trivially_copy_constructible_v<E>}, \tcode{is_trivially_copy_assignable_v<E>}, and \tcode{is_trivially_destructible_v<E>} are all \tcode{true}.
\end{itemdescr}

\indexlibrarymember{operator=}{expected}%
\begin{itemdecl}
constexpr expected& operator=(expected&& rhs) noexcept(is_nothrow_move_constructible_v<E> &&
                                                       is_nothrow_move_assignable_v<E>);
\end{itemdecl}

\begin{itemdescr}
\pnum
\constraints
\tcode{is_reference_v<E> || (is_move_constructible_v<E> \&\& is_move_assignable_v<E>)} is \tcode{true} and \tcode{(is_trivially_move_constructible_v<E> \&\& is_trivially_move_assignable_v<E> \&\& is_trivially_destructible_v<E>)} is \tcode{false}.

\pnum
\effects
If \tcode{this->has_value() \&\& rhs.has_value()}, no effects. Otherwise, if \tcode{this->has_value()}, equivalent to: \tcode{construct_at(addressof(\exposid{unex}), std::move(rhs.\exposid{unex})); \exposid{has_val} = false;} Otherwise, if \tcode{rhs.has_value()}, destroys \exposid{unex} and sets \exposid{has_val} to \tcode{true}. Otherwise, equivalent to \tcode{\exposid{unex} = std::move(rhs.\exposid{unex})}.

\pnum
\returns
\tcode{*this}.

\pnum
\remarks
The exception specification is equivalent to \tcode{is_nothrow_move_constructible_v<E> \&\& is_nothrow_move_assignable_v<E>}. This operator is trivial if \tcode{is_trivially_move_constructible_v<E>}, \tcode{is_trivially_move_assignable_v<E>}, and \tcode{is_trivially_destructible_v<E>} are all \tcode{true}.
\end{itemdescr}

\indexlibrarymember{operator=}{expected}%
\begin{itemdecl}
template<class G>
  constexpr expected& operator=(const unexpected<G>& e);
template<class G>
  constexpr expected& operator=(unexpected<G>&& e);
\end{itemdecl}

\begin{itemdescr}
\pnum
\constraints
\tcode{is_constructible_v<E, const G\&>} is \tcode{true} and \tcode{is_assignable_v<E\&, const G\&>} is \tcode{true}.

\pnum
\effects
If \tcode{has_value()} is \tcode{true}, equivalent to: \tcode{construct_at(addressof(\exposid{unex}), e.error()); \exposid{has_val} = false;} Otherwise, equivalent to: \tcode{\exposid{unex} = unexpected<E>(e.error());}

\pnum
\returns
\tcode{*this}.
\end{itemdescr}

\indexlibrarymember{operator=}{expected}%
\begin{itemdecl}
template<class G>
  constexpr expected& operator=(const unexpected<G>& e);
template<class G>
  constexpr expected& operator=(unexpected<G>&& e);
\end{itemdecl}

\begin{itemdescr}
\pnum
\constraints
\tcode{is_reference_v<G>} is \tcode{true}; and \tcode{is_constructible_v<E, G>} is \tcode{true}; and \tcode{reference_constructs_from_temporary_v<E, G>} is \tcode{false}.

\pnum
\effects
Rebinds \exposid{unex} to refer to the same object as \tcode{e.error()}.

\pnum
\ensures
\tcode{has_value()} is \tcode{false}.

\pnum
\returns
\tcode{*this}.

\pnum
\remarks
This operator never throws: the referent is bound, not copied.
\end{itemdescr}

\indexlibrarymember{operator=}{expected}%
\begin{itemdecl}
template<class G>
  constexpr expected& operator=(const unexpected<G>&) = delete;
\end{itemdecl}

\begin{itemdescr}
\pnum
\constraints
\tcode{is_reference_v<E>} is \tcode{true} and \tcode{is_reference_v<G>} is \tcode{false}.

\pnum
\remarks
When \tcode{E} is a reference type, an overload taking \tcode{unexpected<G>} for a non-reference \tcode{G} is defined as deleted: it would rebind \exposid{unex} to \tcode{unexpected<G>}'s temporary storage.
\end{itemdescr}

\indexlibrarymember{emplace}{expected}%
\begin{itemdecl}
constexpr void emplace() noexcept;
\end{itemdecl}

\begin{itemdescr}
\pnum
\effects
If \tcode{has_value()} is \tcode{false}, destroys \exposid{unex} and sets \exposid{has_val} to \tcode{true}.
\end{itemdescr}

\rSec3[expected.void.swap]{Swap}

\indexlibrarymember{swap}{expected}%
\begin{itemdecl}
constexpr void swap(expected& rhs) noexcept(is_nothrow_move_constructible_v<E> &&
                                            (is_reference_v<E> || is_nothrow_swappable_v<E>));
\end{itemdecl}

\begin{itemdescr}
\pnum
\constraints
\tcode{is_reference_v<E> || is_swappable_v<E>} is \tcode{true} and \tcode{is_move_constructible_v<E>} is \tcode{true}.

\pnum
\effects
If \tcode{this->has_value()} and \tcode{rhs.has_value()}, no effects. If neither \tcode{*this} nor \tcode{rhs} contains a value, equivalent to \tcode{using std::swap; swap(\exposid{unex}, rhs.\exposid{unex});}. If \tcode{rhs.has_value()} is \tcode{false} and \tcode{this->has_value()} is \tcode{true}, initializes \tcode{rhs.\exposid{unex}} from \tcode{std::move(\exposid{unex})}, destroys \exposid{unex}, and leaves \tcode{has_value()} \tcode{false} and \tcode{rhs.has_value()} \tcode{true}. If \tcode{rhs.has_value()} is \tcode{true} and \tcode{this->has_value()} is \tcode{false}, equivalent to \tcode{rhs.swap(*this)}.

\pnum
\throws
Any exception thrown by the expressions in the \Fundescx{Effects}.

\pnum
\remarks
The exception specification is equivalent to \tcode{is_nothrow_move_constructible_v<E> \&\& (is_reference_v<E> || is_nothrow_swappable_v<E>)}.
\end{itemdescr}

\indexlibraryglobal{swap}%
\begin{itemdecl}
friend constexpr void swap(expected& x, expected& y) noexcept(noexcept(x.swap(y)));
\end{itemdecl}

\begin{itemdescr}
\pnum
\constraints
\tcode{is_reference_v<E> || is_swappable_v<E>} is \tcode{true} and \tcode{is_move_constructible_v<E>} is \tcode{true}.

\pnum
\effects
Equivalent to \tcode{x.swap(y)}.
\end{itemdescr}

\rSec3[expected.void.obs]{Observers}

\indexlibrarymember{operator bool}{expected}%
\indexlibrarymember{has_value}{expected}%
\begin{itemdecl}
constexpr explicit operator bool() const noexcept;
constexpr bool has_value() const noexcept;
\end{itemdecl}

\begin{itemdescr}
\pnum
\returns
\exposid{has_val}.
\end{itemdescr}

\indexlibrarymember{operator*}{expected}%
\begin{itemdecl}
constexpr void operator*() const noexcept;
\end{itemdecl}

\begin{itemdescr}
\pnum
\hardexpects
\tcode{has_value()} is \tcode{true}.
\end{itemdescr}

\indexlibrarymember{value}{expected}%
\begin{itemdecl}
constexpr void value() const&;
\end{itemdecl}

\begin{itemdescr}
\pnum
\mandates
\tcode{is_copy_constructible_v<E>} is \tcode{true}.

\pnum
\throws
\tcode{bad_expected_access(error())} if \tcode{has_value()} is \tcode{false}.
\end{itemdescr}

\indexlibrarymember{value}{expected}%
\begin{itemdecl}
constexpr void value() &&;
\end{itemdecl}

\begin{itemdescr}
\pnum
\mandates
\tcode{is_copy_constructible_v<E>} is \tcode{true} and \tcode{is_move_constructible_v<E>} is \tcode{true}.

\pnum
\throws
\tcode{bad_expected_access(std::move(error()))} if \tcode{has_value()} is \tcode{false}.
\end{itemdescr}

\indexlibrarymember{error}{expected}%
\begin{itemdecl}
constexpr const E& error() const& noexcept;
constexpr E& error() & noexcept;
\end{itemdecl}

\begin{itemdescr}
\pnum
\hardexpects
\tcode{has_value()} is \tcode{false}.

\pnum
\returns
\tcode{\exposid{unex}.error()}.
\end{itemdescr}

\indexlibrarymember{error}{expected}%
\begin{itemdecl}
constexpr const E&& error() const&& noexcept;
constexpr E&& error() && noexcept;
\end{itemdecl}

\begin{itemdescr}
\pnum
\hardexpects
\tcode{has_value()} is \tcode{false}.

\pnum
\returns
\tcode{std::move(\exposid{unex}).error()}.
\end{itemdescr}

\indexlibrarymember{error_or}{expected}%
\begin{itemdecl}
template<class G = @\exposid{error-value-type}@>
  constexpr @\exposid{error-value-type}@ error_or(G&& def) const&;
\end{itemdecl}

\begin{itemdescr}
\pnum
\constraints
\tcode{is_copy_constructible_v<remove_cv_t<remove_reference_t<E>>>} is \tcode{true} and \tcode{is_convertible_v<G, remove_cv_t<remove_reference_t<E>>>} is \tcode{true}.

\pnum
\mandates
\tcode{is_copy_constructible_v<\exposid{error-value-type}>} is \tcode{true} and \tcode{is_convertible_v<G, \exposid{error-value-type}>} is \tcode{true}.

\pnum
\returns
\tcode{std::forward<G>(def)} if \tcode{has_value()} is \tcode{true}, \tcode{error()} otherwise.
\end{itemdescr}

\indexlibrarymember{error_or}{expected}%
\begin{itemdecl}
template<class G = @\exposid{error-value-type}@>
  constexpr @\exposid{error-value-type}@ error_or(G&& def) &&;
\end{itemdecl}

\begin{itemdescr}
\pnum
\constraints
\tcode{is_move_constructible_v<remove_cv_t<remove_reference_t<E>>>} is \tcode{true} and \tcode{is_convertible_v<G, remove_cv_t<remove_reference_t<E>>>} is \tcode{true}.

\pnum
\mandates
\tcode{is_move_constructible_v<\exposid{error-value-type}>} is \tcode{true} and \tcode{is_convertible_v<G, \exposid{error-value-type}>} is \tcode{true}.

\pnum
\returns
\tcode{std::forward<G>(def)} if \tcode{has_value()} is \tcode{true}, \tcode{std::move(error())} otherwise.
\end{itemdescr}

\indexlibrarymember{value_or}{expected}%
\begin{itemdecl}
template<class U>
  constexpr void value_or(U&&) const = delete;
\end{itemdecl}

\begin{itemdescr}
\pnum
\constraints
\tcode{is_reference_v<E>} is \tcode{true}.

\pnum
\remarks
\tcode{expected<void, E>} has no \tcode{value_or} member: there is no value to fall back from. This overload exists only to give a clear diagnostic when \tcode{E} is a reference type, and is defined as deleted.
\end{itemdescr}

\rSec3[expected.void.monadic]{Monadic operations}

\indexlibrarymember{and_then}{expected}%
\begin{itemdecl}
template<class F>
  constexpr auto and_then(F&& f) &;
template<class F>
  constexpr auto and_then(F&& f) const&;
\end{itemdecl}

\begin{itemdescr}
\pnum
\constraints
\tcode{is_constructible_v<E, decltype(error())>} is \tcode{true}.

\pnum
\mandates
\tcode{remove_cvref_t<invoke_result_t<F>>} is a specialization of \tcode{expected} and its \tcode{error_type} is the same type as \tcode{E}.

\pnum
\effects
Equivalent to: \tcode{if (has_value()) return invoke(std::forward<F>(f)); else return U(unexpect, error());} where \tcode{U} is \tcode{remove_cvref_t<invoke_result_t<F>>}.
\end{itemdescr}

\indexlibrarymember{and_then}{expected}%
\begin{itemdecl}
template<class F>
  constexpr auto and_then(F&& f) &&;
template<class F>
  constexpr auto and_then(F&& f) const&&;
\end{itemdecl}

\begin{itemdescr}
\pnum
\constraints
\tcode{is_constructible_v<E, decltype(std::move(error()))>} is \tcode{true}.

\pnum
\mandates
\tcode{remove_cvref_t<invoke_result_t<F>>} is a specialization of \tcode{expected} and its \tcode{error_type} is the same type as \tcode{E}.

\pnum
\effects
Equivalent to: \tcode{if (has_value()) return invoke(std::forward<F>(f)); else return U(unexpect, std::move(error()));} where \tcode{U} is \tcode{remove_cvref_t<invoke_result_t<F>>}.
\end{itemdescr}

\indexlibrarymember{or_else}{expected}%
\begin{itemdecl}
template<class F>
  constexpr auto or_else(F&& f) &;
template<class F>
  constexpr auto or_else(F&& f) const&;
\end{itemdecl}

\begin{itemdescr}
\pnum
\mandates
\tcode{remove_cvref_t<invoke_result_t<F, decltype(error())>>} is a specialization of \tcode{expected} and its \tcode{value_type} is the same type as \tcode{T}.

\pnum
\effects
Equivalent to: \tcode{if (has_value()) return G(); else return invoke(std::forward<F>(f), error());} where \tcode{G} is \tcode{remove_cvref_t<invoke_result_t<F, decltype(error())>>}.
\end{itemdescr}

\indexlibrarymember{or_else}{expected}%
\begin{itemdecl}
template<class F>
  constexpr auto or_else(F&& f) &&;
template<class F>
  constexpr auto or_else(F&& f) const&&;
\end{itemdecl}

\begin{itemdescr}
\pnum
\mandates
\tcode{remove_cvref_t<invoke_result_t<F, decltype(std::move(error()))>>} is a specialization of \tcode{expected} and its \tcode{value_type} is the same type as \tcode{T}.

\pnum
\effects
Equivalent to: \tcode{if (has_value()) return G(); else return invoke(std::forward<F>(f), std::move(error()));} where \tcode{G} is \tcode{remove_cvref_t<invoke_result_t<F, decltype(std::move(error()))>>}.
\end{itemdescr}

\indexlibrarymember{transform}{expected}%
\begin{itemdecl}
template<class F>
  constexpr auto transform(F&& f) &;
template<class F>
  constexpr auto transform(F&& f) const&;
\end{itemdecl}

\begin{itemdescr}
\pnum
\constraints
\tcode{is_constructible_v<E, decltype(error())>} is \tcode{true}.

\pnum
\mandates
\tcode{U} is a valid value type for \tcode{expected}, where \tcode{U} is \tcode{remove_cv_t<invoke_result_t<F>>}.

\pnum
\effects
If \tcode{has_value()} is \tcode{false}, returns \tcode{expected<U, E>(unexpect, error())}. Otherwise, if \tcode{is_void_v<U>} is \tcode{false}, returns an \tcode{expected<U, E>} object whose \exposid{has_val} member is \tcode{true} and \tcode{val} member is direct-non-list-initialized with \tcode{invoke(std::forward<F>(f))}. Otherwise, evaluates \tcode{invoke(std::forward<F>(f))} and then returns \tcode{expected<U, E>()}.
\end{itemdescr}

\indexlibrarymember{transform}{expected}%
\begin{itemdecl}
template<class F>
  constexpr auto transform(F&& f) &&;
template<class F>
  constexpr auto transform(F&& f) const&&;
\end{itemdecl}

\begin{itemdescr}
\pnum
\constraints
\tcode{is_constructible_v<E, decltype(std::move(error()))>} is \tcode{true}.

\pnum
\mandates
\tcode{U} is a valid value type for \tcode{expected}, where \tcode{U} is \tcode{remove_cv_t<invoke_result_t<F>>}.

\pnum
\effects
If \tcode{has_value()} is \tcode{false}, returns \tcode{expected<U, E>(unexpect, std::move(error()))}. Otherwise, if \tcode{is_void_v<U>} is \tcode{false}, returns an \tcode{expected<U, E>} object whose \exposid{has_val} member is \tcode{true} and \tcode{val} member is direct-non-list-initialized with \tcode{invoke(std::forward<F>(f))}. Otherwise, evaluates \tcode{invoke(std::forward<F>(f))} and then returns \tcode{expected<U, E>()}.
\end{itemdescr}

\indexlibrarymember{transform_error}{expected}%
\begin{itemdecl}
template<class F>
  constexpr auto transform_error(F&& f) &;
template<class F>
  constexpr auto transform_error(F&& f) const&;
\end{itemdecl}

\begin{itemdescr}
\pnum
\mandates
\tcode{G} is a valid template argument for \tcode{unexpected} and the declaration \tcode{G g(invoke(std::forward<F>(f), error()));} is well-formed, where \tcode{G} is \tcode{remove_cv_t<invoke_result_t<F, decltype(error())>>}.

\pnum
\returns
If \tcode{has_value()} is \tcode{true}, \tcode{expected<T, G>()}; otherwise, an \tcode{expected<T, G>} object whose \exposid{has_val} member is \tcode{false} and \exposid{unex} member is direct-non-list-initialized with \tcode{invoke(std::forward<F>(f), error())}.
\end{itemdescr}

\indexlibrarymember{transform_error}{expected}%
\begin{itemdecl}
template<class F>
  constexpr auto transform_error(F&& f) &&;
template<class F>
  constexpr auto transform_error(F&& f) const&&;
\end{itemdecl}

\begin{itemdescr}
\pnum
\mandates
\tcode{G} is a valid template argument for \tcode{unexpected} and the declaration \tcode{G g(invoke(std::forward<F>(f), std::move(error())));} is well-formed, where \tcode{G} is \tcode{remove_cv_t<invoke_result_t<F, decltype(std::move(error()))>>}.

\pnum
\returns
If \tcode{has_value()} is \tcode{true}, \tcode{expected<T, G>()}; otherwise, an \tcode{expected<T, G>} object whose \exposid{has_val} member is \tcode{false} and \exposid{unex} member is direct-non-list-initialized with \tcode{invoke(std::forward<F>(f), std::move(error()))}.
\end{itemdescr}

\rSec3[expected.void.eq]{Equality operators}

\indexlibraryglobal{operator==}%
\begin{itemdecl}
template<class T2, class E2>
  friend constexpr bool operator==(const expected& x, const expected<T2, E2>& y);
\end{itemdecl}

\begin{itemdescr}
\pnum
\constraints
\tcode{is_void_v<T2>} is \tcode{true}.

\pnum
\mandates
\tcode{is_void_v<T2>} is \tcode{true}. The expression \tcode{x.error() == y.error()} is well-formed and its result is convertible to \tcode{bool}.

\pnum
\returns
If \tcode{x.has_value() != y.has_value()}, \tcode{false}; otherwise, if \tcode{x.has_value()} is \tcode{true}, \tcode{true}; otherwise \tcode{x.error() == y.error()}.
\end{itemdescr}

\indexlibraryglobal{operator==}%
\begin{itemdecl}
template<class E2>
  friend constexpr bool operator==(const expected& x, const unexpected<E2>& e);
\end{itemdecl}

\begin{itemdescr}
\pnum
\mandates
The expression \tcode{x.error() == e.error()} is well-formed and its result is convertible to \tcode{bool}.

\pnum
\returns
\tcode{!x.has_value() \&\& static_cast<bool>(x.error() == e.error())}.
\end{itemdescr}
